% GENERATED FILE. Edit canonical lessons, then run npm run generate:books.
% canonical-source-hash: fnv1a64:b17ca2bc5b08c282
% canonical-lessons: TE-C436-palgonu, TE-C436-todpadu, TE-C436-kudagattu, TE-C436-akarsincu, TE-C436-vibhajincu

\chapter{Taking Part}
\label{ch:te-taking-part}
\hlchaptermodality{436}

\begin{chapteropening}
\textbf{By the end of this chapter:} \emph{I can talk about taking part, contributing, rallying support and dividing.}

Complete the last lesson of chapter 436: I can talk about taking part, contributing, rallying support and dividing.
\end{chapteropening}

\section[\texorpdfstring{\te{పాల్గొను} (pālgonu)}{పాల్గొను (pālgonu)}]{\te{పాల్గొను} (pālgonu) — to take part; to participate}

\label{lesson:TE-C436-palgonu}



\begin{quote}
Before the new one: say the Telugu for a reformer, then the Telugu for to inspire.
\end{quote}

\subsection*{You'll want to know: \te{పాల్గొను}}
\textbf{\te{పాల్గొను}} — \emph{pālgonu} — \textquotedblleft{}to take part; to participate\textquotedblright{}.

\begin{quote}
\te{చాలా} \te{మంది} \te{విద్యార్థులు} \te{ఆ} \te{పోటీలో} \te{పాల్గొన్నారు}.
\end{quote}

\begin{quote}
\emph{cālā mandi vidyārthulu ā pōṭīlō pālgonnāru.}
\end{quote}

\begin{quote}
Many students took part in that competition.
\end{quote}

Past (I …): \te{పాల్గొన్నాను}, \emph{pālgonnānu}.

What you take part in takes -\te{లో}.

\begin{cousinweb}
\textbf{\te{పాల్గొను}} (\emph{pālgonu}) is native Dravidian, a compound of \emph{pālu}, \textquotedblleft{}a part, a share, a lot\textquotedblright{} (DEDR 4097), and \emph{konu}, \textquotedblleft{}to take\textquotedblright{} (DEDR 2151): to take one’s share.
\end{cousinweb}

\begin{grammarlens}[title={What you've built}]
One more word for saying what happens if.
\end{grammarlens}

\subsection*{Your turn}
\begin{itemize}
\raggedright
  \item \emph{Say it:} \emph{pālgonu}
  \item \emph{Say it:} \emph{pālgonu}, once more
  \item \emph{Say it:} \emph{cālā mandi vidyārthulu ā pōṭīlō pālgonnāru.} — Many students took part in that competition.
  \item \emph{Recall:} say the Telugu for a reformer, then the Telugu for to inspire, then say \emph{pālgonu} again
\end{itemize}

\begin{tcolorbox}[breakable,colback=teal!4,colframe=teal!35!black,title={Before you move on}]
What does \te{పాల్గొను} mean? (\textquotedblleft{}To take part; to participate\textquotedblright{}.) Say it once more.
\end{tcolorbox}
\section[\texorpdfstring{\te{తోడ్పడు} (tōḍpaḍu)}{తోడ్పడు (tōḍpaḍu)}]{\te{తోడ్పడు} (tōḍpaḍu) — to contribute; to help along}

\label{lesson:TE-C436-todpadu}



\begin{quote}
Before the new one: say the Telugu for to inspire, then the Telugu for to take part.
\end{quote}

\subsection*{You'll want to know: \te{తోడ్పడు}}
\textbf{\te{తోడ్పడు}} — \emph{tōḍpaḍu} — \textquotedblleft{}to contribute; to help along\textquotedblright{}.

\begin{quote}
\te{ఆ} \te{పథకం} \te{గ్రామ} \te{అభివృద్ధికి} \te{తోడ్పడింది}.
\end{quote}

\begin{quote}
\emph{ā pathakaṁ grāma abhivṛddhiki tōḍpaḍindi.}
\end{quote}

\begin{quote}
That scheme helped the village develop.
\end{quote}

Past (I …): \te{తోడ్పడ్డాను}, \emph{tōḍpaḍḍānu}.

What it helps takes -\te{కి}. The noun is \te{తోడ్పాటు}, a contribution.

\begin{cousinweb}
\textbf{\te{తోడ్పడు}} (\emph{tōḍpaḍu}) is native Telugu, a compound of \emph{tōḍu}, \textquotedblleft{}a companion, help\textquotedblright{}, and the verb \emph{paḍu} (DEDR 3853).
\end{cousinweb}

\begin{grammarlens}[title={What you've built}]
One more word for saying what happens if.
\end{grammarlens}

\subsection*{Your turn}
\begin{itemize}
\raggedright
  \item \emph{Say it:} \emph{tōḍpaḍu}
  \item \emph{Say it:} \emph{tōḍpaḍu}, once more
  \item \emph{Say it:} \emph{ā pathakaṁ grāma abhivṛddhiki tōḍpaḍindi.} — That scheme helped the village develop.
  \item \emph{Recall:} say the Telugu for to inspire, then the Telugu for to take part, then say \emph{tōḍpaḍu} again
\end{itemize}

\begin{tcolorbox}[breakable,colback=teal!4,colframe=teal!35!black,title={Before you move on}]
What does \te{తోడ్పడు} mean? (\textquotedblleft{}To contribute; to help along\textquotedblright{}.) Say it once more.
\end{tcolorbox}
\section[\texorpdfstring{\te{కూడగట్టు} (kūḍagaṭṭu)}{కూడగట్టు (kūḍagaṭṭu)}]{\te{కూడగట్టు} (kūḍagaṭṭu) — to rally; to gather (support)}

\label{lesson:TE-C436-kudagattu}



\begin{quote}
Before the new one: say the Telugu for to take part, then the Telugu for to contribute.
\end{quote}

\subsection*{You'll want to know: \te{కూడగట్టు}}
\textbf{\te{కూడగట్టు}} — \emph{kūḍagaṭṭu} — \textquotedblleft{}to rally; to gather (support)\textquotedblright{}.

\begin{quote}
\te{ఆమె} \te{గ్రామస్తుల} \te{మద్దతు} \te{కూడగట్టింది}.
\end{quote}

\begin{quote}
\emph{āme grāmastula maddatu kūḍagaṭṭindi.}
\end{quote}

\begin{quote}
She rallied the villagers’ support.
\end{quote}

Past (I …): \te{కూడగట్టాను}, \emph{kūḍagaṭṭānu}.

\te{కూడ}, together, + \te{కట్టు}, to tie, sounded \te{గట్టు} after it: to tie people together.

\begin{cousinweb}
\textbf{\te{కూడగట్టు}} (\emph{kūḍagaṭṭu}) is native Dravidian, a compound of \emph{kūḍu}, \textquotedblleft{}to join, to gather, to collect\textquotedblright{} (DEDR 1882), and \emph{kaṭṭu}, \textquotedblleft{}to tie, to bind\textquotedblright{} (DEDR 1147), whose k sounds as g after the first word.
\end{cousinweb}

\begin{grammarlens}[title={What you've built}]
One more word for saying what happens if.
\end{grammarlens}

\subsection*{Your turn}
\begin{itemize}
\raggedright
  \item \emph{Say it:} \emph{kūḍagaṭṭu}
  \item \emph{Say it:} \emph{kūḍagaṭṭu}, once more
  \item \emph{Say it:} \emph{āme grāmastula maddatu kūḍagaṭṭindi.} — She rallied the villagers’ support.
  \item \emph{Recall:} say the Telugu for to take part, then the Telugu for to contribute, then say \emph{kūḍagaṭṭu} again
\end{itemize}

\begin{tcolorbox}[breakable,colback=teal!4,colframe=teal!35!black,title={Before you move on}]
What does \te{కూడగట్టు} mean? (\textquotedblleft{}To rally; to gather (support)\textquotedblright{}.) Say it once more.
\end{tcolorbox}
\section[\texorpdfstring{\te{ఆకర్షించు} (ākarṣiñcu)}{ఆకర్షించు (ākarṣiñcu)}]{\te{ఆకర్షించు} (ākarṣiñcu) — to attract}

\label{lesson:TE-C436-akarsincu}



\begin{quote}
Before the new one: say the Telugu for to contribute, then the Telugu for to rally.
\end{quote}

\subsection*{You'll want to know: \te{ఆకర్షించు}}
\textbf{\te{ఆకర్షించు}} — \emph{ākarṣiñcu} — \textquotedblleft{}to attract\textquotedblright{}.

\begin{quote}
\te{ఆ} \te{కొత్త} \te{పథకం} \te{చాలా} \te{మందిని} \te{ఆకర్షించింది}.
\end{quote}

\begin{quote}
\emph{ā kotta pathakaṁ cālā mandini ākarṣiñcindi.}
\end{quote}

\begin{quote}
That new scheme attracted many people.
\end{quote}

Past (I …): \te{ఆకర్షించాను}, \emph{ākarṣiñcānu}.

The noun is \te{ఆకర్షణ}, attraction.

\begin{cousinweb}
\textbf{\te{ఆకర్షించు}} (\emph{ākarṣiñcu}) is built on Sanskrit \emph{ākarṣaṇa}, \textquotedblleft{}pulling, drawing near, attracting\textquotedblright{} (Monier-Williams), with the Telugu verb ending -\te{ఇంచు}.
\end{cousinweb}

\begin{grammarlens}[title={What you've built}]
One more word for saying what happens if.
\end{grammarlens}

\subsection*{Your turn}
\begin{itemize}
\raggedright
  \item \emph{Say it:} \emph{ākarṣiñcu}
  \item \emph{Say it:} \emph{ākarṣiñcu}, once more
  \item \emph{Say it:} \emph{ā kotta pathakaṁ cālā mandini ākarṣiñcindi.} — That new scheme attracted many people.
  \item \emph{Recall:} say the Telugu for to contribute, then the Telugu for to rally, then say \emph{ākarṣiñcu} again
\end{itemize}

\begin{tcolorbox}[breakable,colback=teal!4,colframe=teal!35!black,title={Before you move on}]
What does \te{ఆకర్షించు} mean? (\textquotedblleft{}To attract\textquotedblright{}.) Say it once more.
\end{tcolorbox}
\section[\texorpdfstring{\te{విభజించు} (vibhajiñcu)}{విభజించు (vibhajiñcu)}]{\te{విభజించు} (vibhajiñcu) — to divide; to split}

\label{lesson:TE-C436-vibhajincu}



\begin{quote}
Before the new one: say the Telugu for to rally, then the Telugu for to attract.
\end{quote}

\subsection*{You'll want to know: \te{విభజించు}}
\textbf{\te{విభజించు}} — \emph{vibhajiñcu} — \textquotedblleft{}to divide; to split\textquotedblright{}.

\begin{quote}
\te{అమ్మ} \te{రొట్టెను} \te{నాలుగు} \te{ముక్కలుగా} \te{విభజించింది}.
\end{quote}

\begin{quote}
\emph{amma roṭṭenu nālugu mukkalugā vibhajiñcindi.}
\end{quote}

\begin{quote}
Mother divided the bread into four pieces.
\end{quote}

Past (I …): \te{విభజించాను}, \emph{vibhajiñcānu}.

The noun is \te{విభజన}, a division.

\begin{cousinweb}
\textbf{\te{విభజించు}} (\emph{vibhajiñcu}) is built on Sanskrit \emph{vibhajana}, \textquotedblleft{}separation, distinction\textquotedblright{} (Monier-Williams), with the Telugu verb ending -\te{ఇంచు}.
\end{cousinweb}

\begin{grammarlens}[title={What you've built}]
One more word for saying what happens if.
\end{grammarlens}

\subsection*{Your turn}
\begin{itemize}
\raggedright
  \item \emph{Say it:} \emph{vibhajiñcu}
  \item \emph{Say it:} \emph{vibhajiñcu}, once more
  \item \emph{Say it:} \emph{amma roṭṭenu nālugu mukkalugā vibhajiñcindi.} — Mother divided the bread into four pieces.
  \item \emph{Recall:} say the Telugu for to rally, then the Telugu for to attract, then say \emph{vibhajiñcu} again
\end{itemize}

\begin{tcolorbox}[breakable,colback=teal!4,colframe=teal!35!black,title={Before you move on}]
What does \te{విభజించు} mean? (\textquotedblleft{}To divide; to split\textquotedblright{}.) Say it once more.
\end{tcolorbox}
